%% 11-integration.tex -- generated by tools/build_from_docx.py; edit freely

\hypertarget{integration}{%
\section{Integration}\label{integration}}

The claim this framework makes is that any partner contribution can be
admitted, bounded and assured through the same interfaces. A claim of
that kind is only worth making if it is tested against the contributions
that actually exist, including the ones that do not fit neatly, because
a platform that accommodates only the convenient contributions is not a
platform.

This chapter tests the claim against the application catalog and the
algorithms contributed by the work packages, and specifies the seams at
which the framework meets the rest of the architecture.

\hypertarget{four-contribution-classes}{%
\subsection{Four Contribution Classes}\label{four-contribution-classes}}

Every partner contribution enters as exactly one of the four
contribution classes of Section 3.4.4, declared when it is registered in
the catalog.

A role implementation supplies Monitoring System, Analytics Engine,
Decision Engine or Actuator behavior for an agent. This is the ordinary
case, and it covers the change scheduler, failure prediction, wireless
transport reconfiguration, energy-aware load balancing, multi-radio
time-sensitive networking, terrestrial and non-terrestrial traffic
engineering, probabilistic forecasting, key performance indicator
monitoring, traffic steering and radio access technology monitoring,
together with the routing and link-estimation algorithms contributed
from the work package on network intelligence.

A platform service supplies a capability the platform itself consumes
and is not an agent. Conflict management and trust modeling are the two.

A reused capability is an existing service bound rather than deployed:
monitoring stacks, model registries, controllers.

A supporting application operates alongside agents without holding a
role in one, such as the extended reality application and semantic
domain switching taken as a service-level function.

\begin{wptable}{11-1}

\begin{longtable}[]{@{}
  >{\raggedright\arraybackslash}p{(\columnwidth - 4\tabcolsep) * \real{0.2125}}
  >{\raggedright\arraybackslash}p{(\columnwidth - 4\tabcolsep) * \real{0.3499}}
  >{\raggedright\arraybackslash}p{(\columnwidth - 4\tabcolsep) * \real{0.4376}}@{}}
\caption{Contribution classes and the contract each is held
to}\tabularnewline
\toprule\noalign{}
\begin{minipage}[b]{\linewidth}\raggedright
\textbf{Class}
\end{minipage} & \begin{minipage}[b]{\linewidth}\raggedright
\textbf{What it supplies}
\end{minipage} & \begin{minipage}[b]{\linewidth}\raggedright
\textbf{Contract it is held to}
\end{minipage} \\
\midrule\noalign{}
\endfirsthead
\toprule\noalign{}
\begin{minipage}[b]{\linewidth}\raggedright
\textbf{Class}
\end{minipage} & \begin{minipage}[b]{\linewidth}\raggedright
\textbf{What it supplies}
\end{minipage} & \begin{minipage}[b]{\linewidth}\raggedright
\textbf{Contract it is held to}
\end{minipage} \\
\midrule\noalign{}
\endhead
\bottomrule\noalign{}
\endlastfoot
\textbf{Role implementation} & Supplies MS, AE, DE or ACT behavior for
an agent. & The full role contract of the relevant chapter, plus
admission and the activation gate. \\
\textbf{Platform service} & Supplies a capability the platform itself
consumes. Not an agent. & Its own service contract. Holds no control
scope; not subject to the self-approval prohibition. \\
\textbf{Reused capability} & An existing service bound rather than
deployed. & Binding contract: pinned interface profile and compatibility
range. \\
\textbf{Supporting application} & Operates alongside agents without
holding a role in one. & Service-level contract. No agent authority, no
action space. \\
\end{longtable}

\end{wptable}

Table 11-1 states the four classes and the contract each is held to.
Only a role implementation is described by an agent descriptor, admitted
as an agent and held to the full role contracts; the other three are
registered with their service contracts, and appear in a descriptor only
as providers of delegated capabilities or as endpoints of bindings.
Nothing enters the platform outside these four classes. Without this
classification, a contribution such as the conflict manager would either
be forced into a role it does not hold, or would bypass admission
entirely, and both are wrong.

\hypertarget{conflict-management-is-a-platform-service-never-an-agent}{%
\subsubsection{Conflict management is a platform service, never an agent}\label{conflict-management-is-a-platform-service-never-an-agent}}

The conflict manager supplies the fabric that mediates the decisions of
other agents. It therefore does not submit to itself, holds no control
scope over network resources, and is not subject to the prohibition on
self-approval, which does not apply to a component that is not seeking
authority.

It is bound at the activation gate as the mediation endpoint, and its
availability is a precondition of other agents being granted decision
emission rather than the reverse. Its own correctness is assured by the
contributing work package and by the platform's conflict metrics, not by
agent-level guardrails.

A mediator obliged to mediate its own verdicts would not terminate. The
distinction between the thing that grants authority and the things that
receive it has to be structural rather than a matter of configuration.

\hypertarget{trust-modeling-is-consumed-at-admission-and-registration}{%
\subsubsection{Trust modeling is consumed at admission and registration}\label{trust-modeling-is-consumed-at-admission-and-registration}}

Distributed-ledger trust between the management functions maps onto
mechanisms the framework already requires: signed capability
advertisements, artifact provenance, registration credentials and audit
records.

Where present, trust scores and attestations are an admission and
registration input. They may gate whether a domain or an artifact is
accepted, and they enrich the audit chain. They are not a runtime input
to individual decisions. Placing a trust evaluation in the per-decision
path would add a cost to every cycle for a property that changes on the
timescale of registration rather than of decisions; keeping it at
admission preserves the loop budget and still obtains the
accountability.

Continuous trust, where a trust controller lowers an agent's trust at
runtime, is not a per-decision query either. It acts on the agent as an
orchestration-initiated authority change, through the same suspension
mechanism as a service-level-agreement-driven migration. Trust enters at
admission; continuous trust is enforced on the agent from outside.
Neither places a trust call on the fast path.

\hypertarget{generative-and-non-deterministic-components}{%
\subsection{Generative and Non-Deterministic Components}\label{generative-and-non-deterministic-components}}

A generative component that converts user intent into low-level
component interfaces is a real contribution in the catalog, and the
framework's position on it follows from the guardrails rather than from
a judgment about the technology.

As a Monitoring System aid or query interface it is admissible normally.
It produces evidence or human-facing views, and the prohibition on
fabricated measurement holds unchanged: generated content must remain
traceable to collected data.

As an Analytics Engine it is admissible with declared uncertainty and
validity predicates, exactly as any other model.

As a Decision Engine it is admissible only where its output is
constrained to the registered action space with machine-checkable
bounds. Structured rationale and last-mile re-validation apply
unchanged, and the advisory entry pattern of Chapter 9 is the default
starting authority.

As an Actuator it is not admissible as a command generator. Translation
to controller-native commands must be deterministic, because the last
mile has to re-validate a command against bounds, and a free-form
command cannot be bounded in the general case.

\begin{wptable}{11-2}

\begin{longtable}[]{@{}
  >{\raggedright\arraybackslash}p{(\columnwidth - 4\tabcolsep) * \real{0.2000}}
  >{\raggedright\arraybackslash}p{(\columnwidth - 4\tabcolsep) * \real{0.2375}}
  >{\raggedright\arraybackslash}p{(\columnwidth - 4\tabcolsep) * \real{0.5625}}@{}}
\caption{Admissibility of a generative component by role}\tabularnewline
\toprule\noalign{}
\begin{minipage}[b]{\linewidth}\raggedright
\textbf{Role}
\end{minipage} & \begin{minipage}[b]{\linewidth}\raggedright
\textbf{Admissibility}
\end{minipage} & \begin{minipage}[b]{\linewidth}\raggedright
\textbf{Condition}
\end{minipage} \\
\midrule\noalign{}
\endfirsthead
\toprule\noalign{}
\begin{minipage}[b]{\linewidth}\raggedright
\textbf{Role}
\end{minipage} & \begin{minipage}[b]{\linewidth}\raggedright
\textbf{Admissibility}
\end{minipage} & \begin{minipage}[b]{\linewidth}\raggedright
\textbf{Condition}
\end{minipage} \\
\midrule\noalign{}
\endhead
\bottomrule\noalign{}
\endlastfoot
\textbf{MS aid or query interface} & Admissible normally. & Generated
content must remain traceable to collected data; fabricated measurement
remains prohibited. \\
\textbf{AE} & Admissible. & Declared uncertainty and machine-evaluable
validity predicates, as for any model. \\
\textbf{DE} & Admissible, advisory by default. & Output constrained to
the registered action space with machine-checkable bounds; structured
rationale and last-mile re-validation unchanged. \\
\textbf{ACT} & Not admissible as a command generator. & Translation to
controller-native commands must be deterministic, since a free-form
command cannot be bounded in the general case. \\
\end{longtable}

\end{wptable}

Table 11-2 states admissibility by role. The reasoning is uniform. The
framework's guardrails are indifferent to how a decision was produced;
they check what it asks for against declared bounds. That indifference
is exactly what allows a generative component to be admitted safely
where its output space is bounded, and exactly why it must be refused
where the output is itself the command.

\hypertarget{the-eight-seams}{%
\subsection{The Eight Seams}\label{the-eight-seams}}

Eight seams join this framework to the rest of the architecture. Each
has an owner and a closed contract. Table 11-3 states them with their
owners.

The first is the agent service tree: the descriptor declares, per tuple,
the services that tuple requires in order to perform its capability,
together with the bindings that interconnect the tuples into a closed
loop. The descriptor declares the agent's scope as local or global, and
that scope determines the orchestrator the descriptor is issued to. The
resource-facing services of the service orchestration pipeline carry no
agent content and are realized independently.

The second is tuple instantiation at the domain and the return of
resolved configuration. The domain service adapter maps the
domain-service block onto the radio, core or transport controller, and
the DMO instantiates the tuples and the services their portion of the
service tree declares, returning the agent instance identifier, the
endpoints and a realization summary.

The third is bus and binding realization: a binding class per edge, with
binding records carrying measured status.

The fourth is the path from the Decision Engine to the conflict fabric,
as specified in Chapter 10.

The fifth is the path from the Analytics Engine to model logistics and
to the evaluator. Models are resolved through Data Continuum bindings by
digest, carrying input and output schema, feature-set version and
approved-use metadata. The evaluator consumes declared evidence and
recommends; the domain or inter-domain management function enacts.

The sixth is the mapping from application to tuple, which is the
classification of Section 11.1 together with the role contracts.

The seventh is controller and actuation proof: read-back independent of
acknowledgement, with actions lacking a proven verification path
declared unverifiable and authority-capped.

The eighth is northbound exposure. Service application interfaces expose
service-level objectives and assurance to consumers, and agents are
never exposed directly.

\begin{wptable}{11-3}

\begin{longtable}[]{@{}
  >{\raggedright\arraybackslash}p{(\columnwidth - 6\tabcolsep) * \real{0.0562}}
  >{\raggedright\arraybackslash}p{(\columnwidth - 6\tabcolsep) * \real{0.2250}}
  >{\raggedright\arraybackslash}p{(\columnwidth - 6\tabcolsep) * \real{0.5312}}
  >{\raggedright\arraybackslash}p{(\columnwidth - 6\tabcolsep) * \real{0.1875}}@{}}
\caption{The eight integration seams and their owners}\tabularnewline
\toprule\noalign{}
\begin{minipage}[b]{\linewidth}\raggedright
\textbf{\#}
\end{minipage} & \begin{minipage}[b]{\linewidth}\raggedright
\textbf{Seam}
\end{minipage} & \begin{minipage}[b]{\linewidth}\raggedright
\textbf{Contract}
\end{minipage} & \begin{minipage}[b]{\linewidth}\raggedright
\textbf{Owner}
\end{minipage} \\
\midrule\noalign{}
\endfirsthead
\toprule\noalign{}
\begin{minipage}[b]{\linewidth}\raggedright
\textbf{\#}
\end{minipage} & \begin{minipage}[b]{\linewidth}\raggedright
\textbf{Seam}
\end{minipage} & \begin{minipage}[b]{\linewidth}\raggedright
\textbf{Contract}
\end{minipage} & \begin{minipage}[b]{\linewidth}\raggedright
\textbf{Owner}
\end{minipage} \\
\midrule\noalign{}
\endhead
\bottomrule\noalign{}
\endlastfoot
\textbf{1} & Agent service tree & Declared in the descriptor: per-tuple
service requirements, tuple bindings, and the agent scope that selects
IDMO or DMO as the receiving orchestrator. & T2.2 and T2.3 \\
\textbf{2} & Tuple instantiation and resolved config & Domain service
adapter maps the domain-service block to the controller; the DMO
instantiates the tuples and their declared services from catalog
artifacts, returning identifiers, endpoints and a realization summary. &
DMO owners \\
\textbf{3} & Bus and binding realization & Binding class per edge, with
binding records carrying measured status. & DMO and bus owners \\
\textbf{4} & DE to conflict fabric & Register, submit, verdict, notify;
screening; bounded approval with stored validity; token checked at ACT.
& WP4 and T2.3 \\
\textbf{5} & AE to model logistics and evaluator & Models resolved by
digest with schema, feature-set version and approved-use metadata;
evaluator recommends, DMO or IDMO enacts. & WP2, WP5 and T2.3 \\
\textbf{6} & Application to tuple mapping & The contribution classes of
Section 11.1 together with the role contracts. & T2.3 and WP5 \\
\textbf{7} & Controller and actuation proof & Read-back independent of
acknowledgement; actions without a proven verification path declared
unverifiable and authority-capped. & Domain controller owners \\
\textbf{8} & Northbound exposure & Service APIs expose service-level
objectives and assurance; agents are never exposed directly. & WP5 and
intent layer \\
\end{longtable}

\end{wptable}

\hypertarget{capability-and-interface-never-implementation-name}{%
\subsection{Capability and Interface, Never Implementation Name}\label{capability-and-interface-never-implementation-name}}

Functions are selected by capability and interface contract rather than
by fixed implementation name. This is not a stylistic preference; it is
the concrete content of the platform claim.

A partner may replace one routing algorithm with another, or one
forecasting model with another, without any change to the service tree,
the bindings or the guardrails, provided the declared capability and
interface contract are satisfied. This is why substitution does not
require re-admission of the whole agent, and it is the difference
between a platform and an integration.

Admission resolves capability requirements against the domain catalog,
and the binding contract fixes the semantic type and schema version
rather than the vendor or the image.

\hypertarget{admission-at-the-authority-the-evidence-supports}{%
\subsection{Admission at the Authority the Evidence Supports}\label{admission-at-the-authority-the-evidence-supports}}

A contribution that cannot satisfy a role contract in full is admitted
at the authority its evidence supports. It is never refused outright
merely for falling short, and never silently exempted from the contract
it fails.

The spectrum runs, in decreasing authority, from
closed\_\allowbreak{}loop\_\allowbreak{}full, through
closed\_\allowbreak{}loop\_\allowbreak{}reduced with a narrowed scope,
tighter bounds and a lower priority band, through advisory, where
intents are published as recommendations and proceed only when approved,
through monitoring\_\allowbreak{}only, to rejection at admission. The
four activation gate outcomes of Chapter 9 implement this directly.
Table 11-4 states the evidence each level requires.

\begin{wptable}{11-4}

\begin{longtable}[]{@{}
  >{\raggedright\arraybackslash}p{(\columnwidth - 4\tabcolsep) * \real{0.1875}}
  >{\raggedright\arraybackslash}p{(\columnwidth - 4\tabcolsep) * \real{0.3125}}
  >{\raggedright\arraybackslash}p{(\columnwidth - 4\tabcolsep) * \real{0.5000}}@{}}
\caption{Authority spectrum and the evidence each level
requires}\tabularnewline
\toprule\noalign{}
\begin{minipage}[b]{\linewidth}\raggedright
\textbf{Authority}
\end{minipage} & \begin{minipage}[b]{\linewidth}\raggedright
\textbf{Meaning}
\end{minipage} & \begin{minipage}[b]{\linewidth}\raggedright
\textbf{Evidence required}
\end{minipage} \\
\midrule\noalign{}
\endfirsthead
\toprule\noalign{}
\begin{minipage}[b]{\linewidth}\raggedright
\textbf{Authority}
\end{minipage} & \begin{minipage}[b]{\linewidth}\raggedright
\textbf{Meaning}
\end{minipage} & \begin{minipage}[b]{\linewidth}\raggedright
\textbf{Evidence required}
\end{minipage} \\
\midrule\noalign{}
\endhead
\bottomrule\noalign{}
\endlastfoot
\textbf{closed\_\allowbreak{}loop\_\allowbreak{}full} & Full declared
authority. & Every gate check satisfied, including a proven verification
path and rollback route for every declared action. \\
\textbf{closed\_\allowbreak{}loop\_\allowbreak{}reduced} & Narrowed
scope, tighter bounds, lower arbitration priority. & Gate satisfied with
capped actions, or sustained conformance not yet demonstrated. \\
\textbf{advisory} & Decisions emitted as recommendations; actuation
requires approval. & Decision logic not declared interpretable, or
verification path unproven for actions the agent would take. \\
\textbf{monitoring\_\allowbreak{}only} & MS and AE active; no decision
emission. & No actuation authority sought, or gate outcome
observation-only. \\
\textbf{Rejected at admission} & Not admitted. & A structural or
semantic admission check failed. \\
\end{longtable}

\end{wptable}

This is what makes the platform genuinely open without making it unsafe.
A partner algorithm that cannot yet prove verified actuation, or cannot
yet express machine-evaluable validity conditions, still has somewhere
to run and a path to promotion, and the network is never exposed to
authority the contribution has not earned. The alternative arrangements
both fail: refusing everything that falls short produces a platform with
one tenant, and exempting contributions that fall short produces a
platform whose guarantees mean nothing.

\hypertarget{requirements-8}{%
\subsection{Requirements}\label{requirements-8}}

\begin{requirement}{Req-T23-IN-01}{MUST}{Every contribution belongs to one contribution class}
\reqfield{Description / rationale}{%
Every partner contribution shall be registered in the catalog as exactly
one contribution class, role\_\allowbreak{}implementation,
platform\_\allowbreak{}service, reused\_\allowbreak{}capability or
supporting\_\allowbreak{}application, and shall be held to the contract
of that class. Only a role implementation shall be described by an agent
descriptor and admitted as an agent.}
\reqfield{Dependency}{%
Req-T23-FW-25}
\reqfield{Traceability (Forward)}{%
\begin{itemize}
\tightlist
\item
  Class declared at catalog registration.
\item
  Role implementations admitted against their descriptor.
\end{itemize}}
\reqfield{UNITY-6G components}{%
\begin{itemize}
\tightlist
\item
  DMO catalog
\item
  Admission (IDMO/DMO)
\end{itemize}}
\reqfield{Linked KPIs}{%
\begin{itemize}
\tightlist
\item
  KPI\#1: Deliver four distributed management entities (MS, AE, DE and
  ACT) deployable in each heterogeneous domain.
\item
  KPI\#41: Reduction of 70\% of deployment time of AI-based
  resource/energy management algorithms.
\end{itemize}}
\end{requirement}

\begin{requirement}{Req-T23-IN-02}{MUST}{A platform service is not an agent}
\reqfield{Description / rationale}{%
A contribution classified as a platform service shall hold no control
scope over network resources, shall not submit intent for mediation and
shall not be subject to the role contracts of Chapters 4 to 7. The
conflict management service shall be bound at the activation gate as the
mediation endpoint, and its availability shall be a precondition of
other agents being granted decision emission.}
\reqfield{Dependency}{%
Req-T23-CG-04}
\reqfield{Traceability (Forward)}{%
\begin{itemize}
\tightlist
\item
  Class checked at admission.
\item
  Mediation endpoint bound at gate check 7.
\end{itemize}}
\reqfield{UNITY-6G components}{%
\begin{itemize}
\tightlist
\item
  DMO
\item
  Conflict fabric
\item
  WP4 resolver
\end{itemize}}
\reqfield{Linked KPIs}{%
\begin{itemize}
\tightlist
\item
  KPI\#1: Deliver four distributed management entities (MS, AE, DE and
  ACT) deployable in each heterogeneous domain.
\item
  KPI\#21: Reduce the performance penalty due to conflicts by 20\% by
  means of DLT, hierarchical algorithms and policy-based conflict
  resolution schemes.
\end{itemize}}
\end{requirement}

\begin{requirement}{Req-T23-IN-03}{MUST}{Trust is consumed at admission, never per decision}
\reqfield{Description / rationale}{%
Where a trust service is present, trust scores and attestations shall be
inputs to admission and registration, and may gate acceptance of a
domain or an artifact. They shall not be a runtime input to individual
decisions. A runtime trust change shall act on the agent as an
orchestration-initiated authority change through the suspension
mechanism, and shall place no trust call on the decision path.}
\reqfield{Dependency}{%
Req-T23-LC-05}
\reqfield{Traceability (Forward)}{%
\begin{itemize}
\tightlist
\item
  Trust evaluated at admission and registration.
\item
  Runtime trust changes applied as authority transitions.
\end{itemize}}
\reqfield{UNITY-6G components}{%
\begin{itemize}
\tightlist
\item
  DMO
\item
  IDMO
\item
  Trust service
\end{itemize}}
\reqfield{Linked KPIs}{%
\begin{itemize}
\tightlist
\item
  KPI\#33: Support more tenants via distributed AI \& DLT.
\item
  KPI\#4: Real-time reconfiguration of IEEE 802.11 network domain in
  range smaller than 10 ms from taking the decision.
\end{itemize}}
\end{requirement}

\begin{requirement}{Req-T23-IN-04}{MUST}{Generative components are bounded by role}
\reqfield{Description / rationale}{%
A generative or non-deterministic component shall be admissible as an MS
aid or query interface with generated content traceable to collected
data; as an AE with declared uncertainty and validity predicates; and as
a DE only where its output is constrained to the registered action space
with machine-checkable bounds, entering at advisory by default.}
\reqfield{Dependency}{%
Req-T23-FW-18}
\reqfield{Traceability (Forward)}{%
\begin{itemize}
\tightlist
\item
  Output bounds checked at admission and re-validated at the last mile.
\end{itemize}}
\reqfield{UNITY-6G components}{%
\begin{itemize}
\tightlist
\item
  Admission (IDMO/DMO)
\item
  DE runtime
\item
  ACT runtime
\end{itemize}}
\reqfield{Linked KPIs}{%
\begin{itemize}
\tightlist
\item
  KPI\#36: Guarantee 99.99999\% service availability and reliability
  based on effective monitoring and LCM.
\item
  KPI\#41: Reduction of 70\% of deployment time of AI-based
  resource/energy management algorithms.
\end{itemize}}
\end{requirement}

\begin{requirement}{Req-T23-IN-05}{MUST NOT}{Generative command generation is not admissible}
\reqfield{Description / rationale}{%
A generative component shall not be admitted as an ACT command
generator. Translation to controller-native commands shall be
deterministic, so that a command can be re-validated against declared
bounds at the last mile.}
\reqfield{Dependency}{%
Req-T23-FW-17}
\reqfield{Traceability (Forward)}{%
\begin{itemize}
\tightlist
\item
  Checked at admission against the declared ACT translation.
\end{itemize}}
\reqfield{UNITY-6G components}{%
\begin{itemize}
\tightlist
\item
  Admission (IDMO/DMO)
\item
  ACT runtime
\end{itemize}}
\reqfield{Linked KPIs}{%
\begin{itemize}
\tightlist
\item
  KPI\#36: Guarantee 99.99999\% service availability and reliability
  based on effective monitoring and LCM.
\item
  KPI\#42: DT-validated AI algorithms that do not break the
  time-sensitive operation when deployed in 99\% of the outputs.
\end{itemize}}
\end{requirement}

\begin{requirement}{Req-T23-IN-06}{MUST}{Selection is by capability and interface contract}
\reqfield{Description / rationale}{%
Functions shall be selected by declared capability and interface
contract rather than by fixed implementation name. Admission shall
resolve capability requirements against the domain catalog, and the
binding contract shall fix semantic type and schema version rather than
vendor or image. Substitution of a conforming implementation shall not
require re-admission of the whole agent.}
\reqfield{Dependency}{%
Req-T23-IC-10}
\reqfield{Traceability (Forward)}{%
\begin{itemize}
\tightlist
\item
  Capability resolved against the catalog at admission.
\end{itemize}}
\reqfield{UNITY-6G components}{%
\begin{itemize}
\tightlist
\item
  Admission (IDMO/DMO)
\item
  DMO catalog
\end{itemize}}
\reqfield{Linked KPIs}{%
\begin{itemize}
\tightlist
\item
  KPI\#41: Reduction of 70\% of deployment time of AI-based
  resource/energy management algorithms.
\item
  KPI\#1: Deliver four distributed management entities (MS, AE, DE and
  ACT) deployable in each heterogeneous domain.
\end{itemize}}
\end{requirement}

\begin{requirement}{Req-T23-IN-07}{MUST}{Admission grants the authority the evidence supports}
\reqfield{Description / rationale}{%
A contribution that cannot satisfy a role contract in full shall be
admitted at the authority its evidence supports, along the spectrum
closed\_\allowbreak{}loop\_\allowbreak{}full,
closed\_\allowbreak{}loop\_\allowbreak{}reduced, advisory,
monitoring\_\allowbreak{}only and rejection. It shall not be refused
solely for falling short of full authority, and shall not be exempted
from the contract it fails.}
\reqfield{Dependency}{%
Req-T23-FW-23}
\reqfield{Traceability (Forward)}{%
\begin{itemize}
\tightlist
\item
  Implemented by the four activation gate outcomes.
\end{itemize}}
\reqfield{UNITY-6G components}{%
\begin{itemize}
\tightlist
\item
  DMO
\item
  Platform runtime
\end{itemize}}
\reqfield{Linked KPIs}{%
\begin{itemize}
\tightlist
\item
  KPI\#41: Reduction of 70\% of deployment time of AI-based
  resource/energy management algorithms.
\item
  KPI\#1: Deliver four distributed management entities (MS, AE, DE and
  ACT) deployable in each heterogeneous domain.
\end{itemize}}
\end{requirement}

\begin{requirement}{Req-T23-IN-08}{MUST}{Agents are not exposed northbound}
\reqfield{Description / rationale}{%
Northbound service interfaces shall expose service-level objectives and
assurance to consumers. Agents, their descriptors and their internal
state shall not be exposed directly through those interfaces.}
\reqfield{Dependency}{%
Req-T23-LC-08}
\reqfield{Traceability (Forward)}{%
\begin{itemize}
\tightlist
\item
  Exposure boundary held at the intent layer.
\end{itemize}}
\reqfield{UNITY-6G components}{%
\begin{itemize}
\tightlist
\item
  Intent layer
\item
  IDMO
\end{itemize}}
\reqfield{Linked KPIs}{%
\begin{itemize}
\tightlist
\item
  KPI\#1: Deliver four distributed management entities (MS, AE, DE and
  ACT) deployable in each heterogeneous domain.
\item
  KPI\#36: Guarantee 99.99999\% service availability and reliability
  based on effective monitoring and LCM.
\end{itemize}}
\end{requirement}

\hypertarget{contribution-to-the-project-kpis-7}{%
\subsection{CONTRIBUTION TO THE PROJECT KPIs}\label{contribution-to-the-project-kpis-7}}

Table 11-5 states what integration contributes to each KPI it bears on,
and Table 11-6 maps the requirements of this chapter onto those KPIs.

\begin{wptable}{11-5}

\begin{longtable}[]{@{}
  >{\raggedright\arraybackslash}p{(\columnwidth - 4\tabcolsep) * \real{0.1499}}
  >{\raggedright\arraybackslash}p{(\columnwidth - 4\tabcolsep) * \real{0.6375}}
  >{\raggedright\arraybackslash}p{(\columnwidth - 4\tabcolsep) * \real{0.2126}}@{}}
\caption{Contribution of integration to the project KPIs}\tabularnewline
\toprule\noalign{}
\begin{minipage}[b]{\linewidth}\raggedright
\textbf{KPI}
\end{minipage} & \begin{minipage}[b]{\linewidth}\raggedright
\textbf{Contribution of integration}
\end{minipage} & \begin{minipage}[b]{\linewidth}\raggedright
\textbf{Requirements}
\end{minipage} \\
\midrule\noalign{}
\endfirsthead
\toprule\noalign{}
\begin{minipage}[b]{\linewidth}\raggedright
\textbf{KPI}
\end{minipage} & \begin{minipage}[b]{\linewidth}\raggedright
\textbf{Contribution of integration}
\end{minipage} & \begin{minipage}[b]{\linewidth}\raggedright
\textbf{Requirements}
\end{minipage} \\
\midrule\noalign{}
\endhead
\bottomrule\noalign{}
\endlastfoot
\textbf{KPI\#1} & Four contribution classes, each held to its own
contract, let heterogeneous partner contributions share the platform
without any of them being special-cased or bypassing registration. &
Req-T23-IN-01, Req-T23-IN-02 \\
\textbf{KPI\#41} & Selection by capability and interface contract means
a partner may substitute an algorithm without re-admission of the whole
agent. & Req-T23-IN-06, Req-T23-IN-07 \\
\textbf{KPI\#21} & Conflict management is a platform service rather than
an agent, which keeps the mediator structurally separate from the agents
it mediates and lets its verdicts terminate. & Req-T23-IN-02 \\
\textbf{KPI\#33} & DLT-based trust consumed at admission and
registration lets more tenants be admitted with an audit chain behind
them, without placing a trust call on the per-decision path. &
Req-T23-IN-03 \\
\textbf{KPI\#4} & Keeping trust and conflict off the fast path preserves
the loop budget for the tightest latency classes. & Req-T23-IN-03 \\
\textbf{KPI\#36} & A generative component is refused precisely where its
output would be the command, and admitted where its output space is
bounded and checkable, so no free-form command reaches a live service. &
Req-T23-IN-04, Req-T23-IN-05 \\
\end{longtable}

\end{wptable}

\begin{wptable}{11-6}\footnotesize

\begin{longtable}[]{@{}
  >{\raggedright\arraybackslash}p{(\columnwidth - 14\tabcolsep) * \real{0.1811}}
  >{\raggedright\arraybackslash}p{(\columnwidth - 14\tabcolsep) * \real{0.1170}}
  >{\raggedright\arraybackslash}p{(\columnwidth - 14\tabcolsep) * \real{0.1170}}
  >{\raggedright\arraybackslash}p{(\columnwidth - 14\tabcolsep) * \real{0.1170}}
  >{\raggedright\arraybackslash}p{(\columnwidth - 14\tabcolsep) * \real{0.1170}}
  >{\raggedright\arraybackslash}p{(\columnwidth - 14\tabcolsep) * \real{0.1170}}
  >{\raggedright\arraybackslash}p{(\columnwidth - 14\tabcolsep) * \real{0.1170}}
  >{\raggedright\arraybackslash}p{(\columnwidth - 14\tabcolsep) * \real{0.1170}}@{}}
\caption{Mapping of integration requirements onto the project
KPIs}\tabularnewline
\toprule\noalign{}
\begin{minipage}[b]{\linewidth}\raggedright
\textbf{KPI\textbackslash{}}

\textbf{Requirement}
\end{minipage} & \begin{minipage}[b]{\linewidth}\raggedright
\textbf{\#1}
\end{minipage} & \begin{minipage}[b]{\linewidth}\raggedright
\textbf{\#4}
\end{minipage} & \begin{minipage}[b]{\linewidth}\raggedright
\textbf{\#21}
\end{minipage} & \begin{minipage}[b]{\linewidth}\raggedright
\textbf{\#33}
\end{minipage} & \begin{minipage}[b]{\linewidth}\raggedright
\textbf{\#36}
\end{minipage} & \begin{minipage}[b]{\linewidth}\raggedright
\textbf{\#41}
\end{minipage} & \begin{minipage}[b]{\linewidth}\raggedright
\textbf{\#42}
\end{minipage} \\
\midrule\noalign{}
\endfirsthead
\toprule\noalign{}
\begin{minipage}[b]{\linewidth}\raggedright
\textbf{KPI\textbackslash{}}

\textbf{Requirement}
\end{minipage} & \begin{minipage}[b]{\linewidth}\raggedright
\textbf{\#1}
\end{minipage} & \begin{minipage}[b]{\linewidth}\raggedright
\textbf{\#4}
\end{minipage} & \begin{minipage}[b]{\linewidth}\raggedright
\textbf{\#21}
\end{minipage} & \begin{minipage}[b]{\linewidth}\raggedright
\textbf{\#33}
\end{minipage} & \begin{minipage}[b]{\linewidth}\raggedright
\textbf{\#36}
\end{minipage} & \begin{minipage}[b]{\linewidth}\raggedright
\textbf{\#41}
\end{minipage} & \begin{minipage}[b]{\linewidth}\raggedright
\textbf{\#42}
\end{minipage} \\
\midrule\noalign{}
\endhead
\bottomrule\noalign{}
\endlastfoot
\textbf{Req-T23-IN-01} & \kpimark{} & & & & & \kpimark{} & \\
\textbf{Req-T23-IN-02} & \kpimark{} & & \kpimark{} & & & & \\
\textbf{Req-T23-IN-03} & & \kpimark{} & & \kpimark{} & & & \\
\textbf{Req-T23-IN-04} & & & & & \kpimark{} & \kpimark{} & \\
\textbf{Req-T23-IN-05} & & & & & \kpimark{} & & \kpimark{} \\
\textbf{Req-T23-IN-06} & \kpimark{} & & & & & \kpimark{} & \\
\textbf{Req-T23-IN-07} & \kpimark{} & & & & & \kpimark{} & \\
\textbf{Req-T23-IN-08} & \kpimark{} & & & & \kpimark{} & & \\
\end{longtable}

\end{wptable}
